% =========================================================================
%  Ejercicio 1h - Modelos de datos
%  Responsable: [nombre]
% =========================================================================

\section*{Ejercicio 1h. Modelos de datos}

% Responder aqui. Contenido minimo:
% - Tabla comparativa entre al menos CUATRO modelos de datos (jerarquico,
%   de red, relacional, orientado a objetos, NoSQL: llave-valor, documento,
%   columnar, grafo), destacando caracteristicas, ventajas, desventajas y
%   casos de uso tipicos. Se puede usar un entorno \texttt{tabular} o
%   \texttt{tabularx}.
% - Argumentar por que el modelo relacional se volvio el mas adoptado.
% - Contextos actuales donde los modelos NoSQL superan al relacional en
%   eficiencia o flexibilidad, con ejemplos.

% Ejemplo de tabla comparativa (editar y completar):
% \begin{table}[H]
%     \centering
%     \caption{Comparativa de modelos de datos.}
%     \begin{tabularx}{\textwidth}{l X X X X}
%         \toprule
%         \textbf{Modelo} & \textbf{Caracteristicas} & \textbf{Ventajas} &
%         \textbf{Desventajas} & \textbf{Casos de uso} \\
%         \midrule
%         Relacional & ... & ... & ... & ... \\
%         Documento & ... & ... & ... & ... \\
%         Columnares & ... & ... & ... & ... \\
%         Grafos & ... & ... & ... & ... \\
%         \bottomrule
%     \end{tabularx}
% \end{table}